% =============================================================================
% intro.tex — Páginas preliminares
% Orden: Dedicatoria → Agradecimientos → Resumen → Abstract
% Ninguna de estas páginas aparece en el índice de contenido.
% Todas usan \paginaprevia{TÍTULO} definido en univalle-perfil.cls.
% =============================================================================

\paginaprevia{DEDICATORIA}

A mi familia, por el apoyo incondicional durante todo este proceso. Por cada vez que confiaron en mí incluso cuando yo dudaba, y por hacer posible que llegara hasta aquí.

A los docentes que no solo enseñaron, sino que se tomaron el tiempo de guiar. Este trabajo también es producto de lo que sembraron.

\paginaprevia{AGRADECIMIENTOS}

A mi familia, por el esfuerzo y la paciencia que pusieron en acompañarme a lo largo de esta etapa.

A los docentes de la Universidad Privada del Valle que orientaron este proyecto, especialmente a quienes dedicaron tiempo más allá del aula para ayudarme a pensar con claridad.

A los estudiantes del Colegio Cardenal Cushing que participaron en las entrevistas y pruebas, cuya experiencia fue la base real de este trabajo.

\paginaprevia{RESUMEN}

El presente proyecto de grado propone el desarrollo de una plataforma adaptativa de aprendizaje dirigida a estudiantes de último año de colegio en Santa Cruz de la Sierra, Bolivia. El problema que motiva el proyecto es que los estudiantes estudian de manera desordenada y pasiva, sin ninguna herramienta que organice el contenido según su nivel individual, ajuste la dificultad de manera continua, verifique que la comprensión sea real y programe revisiones automáticas para evitar el olvido.

La plataforma se fundamenta en el modelo de procesamiento de la información de Atkinson y Shiffrin (1968), que describe el aprendizaje como un flujo de etapas cognitivas con condiciones específicas para funcionar. A partir de ese marco, el sistema organiza el conocimiento en un grafo dinámico generado automáticamente para cualquier dominio, construye un modelo de usuario individual, implementa un motor de dificultad progresiva y un algoritmo de espaciado automático basado en la curva del olvido, y entrega retroalimentación inmediata después de cada tarea.

El entregable principal es una arquitectura base funcional construida en fases incrementales, validada mediante pruebas técnicas y de usuario. La metodología es de desarrollo iterativo con enfoque mixto de investigación aplicada.

\textbf{Palabras clave:} aprendizaje adaptativo, grafo de conocimiento, espaciado automático, modelo de usuario, procesamiento de la información.

\paginaprevia{ABSTRACT}

This undergraduate project proposes the development of an adaptive learning platform aimed at final-year high school students in Santa Cruz de la Sierra, Bolivia. The problem that motivates the project is that students study in a disorganized and passive way, without any tool that organizes content according to their individual level, continuously adjusts difficulty, verifies real comprehension, or schedules automatic reviews to prevent forgetting.

The platform is grounded in the information processing model proposed by Atkinson and Shiffrin (1968), which describes learning as a flow of cognitive stages with specific conditions for functioning. Based on this framework, the system organizes knowledge into a dynamic graph automatically generated for any domain, builds an individual user model, implements a progressive difficulty engine and an automatic spaced repetition algorithm based on the forgetting curve, and delivers immediate feedback after each task.

The main deliverable is a functional base architecture built in incremental phases, validated through technical and user testing. The methodology follows iterative development with a mixed applied research approach.

\textbf{Keywords:} adaptive learning, knowledge graph, spaced repetition, user model, information processing.
